% Texto del resumen (≈200 palabras): pregunta, datos, método y los tres
% hallazgos principales. Solo el texto: el entorno abstract vive en main.tex
% porque en cas-dc es tipo verbatim y no admite \end{abstract} en un \input.
Se estudia cuántos accidentes de tránsito se registrarán y dónde en la Zona
Metropolitana de Monterrey (ZMM). Los datos son los de ATUS (INEGI) de
2019--2024: 379,294 accidentes en 18 municipios, agregados en 2,270 celdas
hexagonales de 500\,m. Se comparan modelos de conteo con un modelo jerárquico
bayesiano con estructura CAR, efecto municipal dinámico y efecto de celda neto
del municipio. Primero, la concentración espacial es persistente: el 5\,\% de
las celdas reúne el 42.8\,\% de los accidentes (Gini de 0.758; $I$ de Moran de
0.654) y el orden de las celdas se mantiene entre 2019 y 2024 (Spearman de
0.83). Segundo, el municipio explica casi todo el residual espacial: el $I$ de Moran
del residual de la Binomial Negativa en 2024 baja de 0.30 a 0.06 al restar la
media municipal, lo
que es compatible con un régimen de reporte distinto por municipio, sin que la
base permita separarlo del riesgo real. Tercero, para 2024 el modelo jerárquico predice
70,810 accidentes contra 71,295 observados y reduce el CRPS (puntuación de
probabilidad ranqueada continua) de 6.49 a 6.10 frente al mejor modelo de
referencia, pero no mejora el PAI (índice de precisión predictiva, 8.44). Los
registros son administrativos y no incluyen exposición, de modo que se estima
frecuencia y no tasa; la validación usa una sola ventana (2024) y los
resultados son asociaciones, no efectos.
